%=======================================================================
%  CHAPTER 6 — MACHINE LEARNING  (6 hrs)
%=======================================================================
\section{Machine Learning}

{\color{sub}\footnotesize 6 hrs \gap$\cdot$\gap asked in \mk{40 of 41 papers}
\gap$\cdot$\gap 6 syllabus sub-topics \gap$\cdot$\gap 436 marks, 12.9\,\% of every mark
ever set \gap$\cdot$\gap the ID3 data tables, the GA generations and the fuzzy
arithmetic are worked in \emph{Ch6 -- Machine Learning -- Solved Problems}.}

\lead{Read this first:}
\begin{itemize}
  \item \mk{Two topics carry this chapter, and both are drawings.}
        \textbf{Genetic algorithm} \tS{21} and \textbf{fuzzy logic} \tS{16}. Between
        them they are asked in \mk{31 of the 41 papers} and they take
        roughly two thirds of the chapter's marks. Learn the GA flowchart and the
        fuzzy-inference block diagram cold; everything else here is a bonus.
  \item \mk{The examiner pairs them with each other.} \yr{82 Ka},
        \yr{\textbf{\texttt{82 Bh}}} and \yr{\texttt{81 Ba}} all ask fuzzy logic and GA in
        the \emph{same} question. Answering one well and the other badly costs the
        whole 8 marks.
  \item \textbf{Third is supervised Vs unsupervised} \tS{12}, always with the rider
        \emph{``clustering falls under which category? Justify, stating algorithms''}.
        Name \textbf{K-means} and \textbf{hierarchical clustering} or the justification
        scores nothing.
  \item \mk{ID3 is the only real arithmetic in the chapter} \tF{6}, and
        it is always the same shape: entropy of the whole set, entropy of each
        attribute's split, gain $=$ difference, pick the largest. Three different data
        tables appear, all worked in the companion.
  \item \mk{The syllabus line says ``fuzzy learning'' and no paper has
        ever asked that.} Every one of the sixteen asks the \textbf{fuzzy inference
        system} $-$ membership functions, rules, fuzzification, defuzzification. \S6.7
        teaches the inference system and says where the two names meet.
  \item Reliable short notes, each a fixed list you can write in four minutes:
        \textbf{GA} \m{3}\m{4}, \textbf{fuzzy logic} \m{3}\m{4},
        \textbf{Boltzmann machine} \m{4}, \textbf{supervised Vs unsupervised} \m{4},
        \textbf{mutation Vs crossover} \m{3}, \textbf{induction Vs deduction} \m{3}.
\end{itemize}

\hr

%-----------------------------------------------------------------------
\T{6.1 What machine learning is, and the learning framework}

\Q{What is machine learning? Explain the learning framework with a block diagram.}{\m{1}\m{6}\m{2+6}}
{\color{sub}\footnotesize \tP{2} \yr{78 Ba \m{2+6}, \textbf{69 Bh} \m{6}}
\gap$\cdot$\gap the \emph{definition} alone opens eleven more papers, always as the
\m{1} or \m{2} lead-in to a GA, fuzzy or Boltzmann question (\yr{80 Ash \m{1}}
asks nothing else) \gap$\cdot$\gap so learn the definition as one sentence and the
framework as four boxes}

\sbs{%
\lead{Definition}
\begin{itemize}
  \item \textbf{Machine learning is the branch of AI in which a program improves its
        performance on a task from experience, without being explicitly programmed for
        that task.}
  \item Tom Mitchell's form, worth quoting: a program \textbf{learns} from experience
        $E$ w.r.t. task $T$ and measure $P$ if \mk{its performance at $T$,
        measured by $P$, improves with $E$}.
  \item In AI terms: the agent \textbf{observes data}, \textbf{builds a model} of it,
        and uses that model as both a hypothesis about the world and a program that
        solves problems.
  \item \textbf{Three things learning always does}, and the three-word answer to
        \emph{``what is learning''}:
  \begin{itemize}
    \item \textbf{Change} $-$ the learner itself is altered.
    \item \textbf{Generalization} $-$ performance improves on \emph{similar} tasks, not
          only on the ones seen.
    \item \textbf{Improvement} $-$ and the design must stop changes that make
          performance worse.
  \end{itemize}
  \item \mk{Why an intelligent agent needs it} (asked as
        \emph{``learning is an essential characteristic of intelligent agents, comment''},
        Ch\,1): a designer cannot enumerate every situation; the environment changes
        after deployment; and some tasks have no known algorithm, only examples.
\end{itemize}}{%
\lead{The learning framework: four components}
\figT{c6_framework.png}
{\color{sub}\footnotesize Environment feeds the learning element; the learning element
writes the knowledge base; the performance element uses it and reports back.}
\begin{itemize}
  \item \textbf{1. Environment} $-$ supplies the information. Judge it on:
  \begin{itemize}
    \item \textbf{Level}: high level $=$ abstract, covers a broad class of problems;
          low level $=$ detailed, covers one problem.
    \item \textbf{Quality}: noise free, reliable, ordered.
  \end{itemize}
  \item \textbf{2. Learning element} $-$ turns that information into knowledge. Rote
        learning, learning by example, by analogy, explanation-based (\S6.2).
  \item \textbf{3. Knowledge base} $-$ stores what was learnt. Must be
        \textbf{expressive}, \textbf{modifiable} and \textbf{extendible} (holds
        meta-knowledge, so its own structure can change).
  \item \textbf{4. Performance element} $-$ uses the knowledge to act, measures how well
        it did, and \mk{feeds the error back} to the learning element.
        That feedback arrow is the whole point of the diagram: draw it.
\end{itemize}}

%-----------------------------------------------------------------------
\T{6.2 Ways of learning, and induction Vs deduction}

\Q{What are the different ways / methods of learning? Describe each.}{\m{8}\m{3+5}}
{\color{sub}\footnotesize \tP{2} \yr{\textbf{\textit{71 Bh}} \m{8}, \texttt{81 Ba} \m{3+5}}
\gap$\cdot$\gap \yr{\texttt{81 Ba}} asks it as the \m{3} half of a GA question, so give
the list, not essays \gap$\cdot$\gap \yr{\textbf{78 Ch}} \m{2} asks the narrower
\emph{``algorithms inspired by biological phenomena and social behaviour''}: genetic
algorithms, neural networks, swarm intelligence (ant colony, particle swarm),
artificial immune systems}

\sbs{%
\lead{The five ways, easiest to hardest}
\begin{itemize}
  \item \textbf{1. Rote learning} $-$ memorise and recall.
  \begin{itemize}
    \item Solved problem $\rightarrow$ store $\rightarrow$ look up next time.
    \item No inference at all: nothing new is derived.
    \item \mk{Only worth it if retrieval is cheaper than recomputing.}
    \item Eg Samuel's checkers player, which stored board positions it had seen.
  \end{itemize}
  \item \textbf{2. Learning by example (inductive)} $-$ specific cases
        $\rightarrow$ general rule. Eg decision trees, ID3 (\S6.5).
  \item \textbf{3. Learning by analogy} $-$ transfer a solution from a known similar
        problem (\S6.4).
  \item \textbf{4. Explanation-based learning (EBL)} $-$ use existing domain knowledge to
        \emph{explain} one example, then generalise the explanation into a rule.
  \begin{itemize}
    \item Learns from \mk{a single example}, because the prior knowledge
          does the work the extra examples would have done.
    \item Eg see one cup once, know what makes it a cup, recognise every cup after.
  \end{itemize}
  \item \textbf{5. Learning by deduction} $-$ general rules $\rightarrow$ specific facts,
        or new general rules from old ones. $A=B$, $B=C$, so $A=C$.
\end{itemize}}{%
\lead{Induction Vs deduction}
{\color{sub}\footnotesize asked as a short note by \yr{\textbf{\textit{73 Bh}}} \m{3}.}
\figT{c6_ind_ded.png}
{\color{sub}\footnotesize Read the left ladder downwards and the right one upwards:
that direction is the whole distinction.}}

\sbsr{0.62}{%
\begin{tabularx}{\linewidth}{@{}L{2.6cm}XX@{}}
\hrow \thead{} & \thead{Inductive} & \thead{Deductive} \\
Direction & observation $\rightarrow$ hypothesis & hypothesis $\rightarrow$ observation \\
Goes from & specific & general \\
Produces & a general rule from raw data & a specific fact from known rules \\
Truth & \mk{not guaranteed}, only probable & guaranteed if the premises hold \\
Reasoning & probabilistic & logical (Ch\,4) \\
Example & 1000 white swans $\rightarrow$ ``swans are white'' & all men are mortal;
          Socrates is a man; so Socrates is mortal \\
\end{tabularx}}{%
\begin{itemize}
  \item \mk{The one-line answer}: induction generalises and can be
        wrong; deduction specialises and cannot.
  \item Both appear in \S6.2's list of ways of learning, and the papers ask them as a
        \m{3} short note, so three rows of that table plus this line is a full answer.
  \item \mk{Where each one lives in this subject}: deduction is
        Chapter 4 (resolution, forward and backward chaining, all truth-preserving);
        induction is this chapter (ID3, \S6.5) and every learning algorithm in it.
  \item \textbf{Abduction} is worth one extra line if you have room: reasoning from an
        observation to its \emph{most likely cause}. Medical diagnosis. Also not
        truth-preserving.
\end{itemize}}

%-----------------------------------------------------------------------
\T{6.3 Supervised, unsupervised and reinforcement learning}

\Q{Differentiate supervised and unsupervised learning. Clustering falls under which category? Justify, stating algorithms.}{\m{4}\m{5}\m{2+4}\m{4+4}\m{2+6}}
{\color{sub}\footnotesize \tS{12} \yr{\textit{74 Ma} \m{2+4},
\textbf{\textit{72 Ash}} \m{2+4}, \textbf{\textit{74 Bh}} \m{4+4}, 79 Jth \m{4+4},
\textbf{\textit{76 Bh}} \m{5}, 78 Po \m{2+6}, 69 Po \m{4}, \textbf{74 Bh} \m{4},
78 Ba \m{4}, \textbf{\textit{77 Ch}} \m{4}, 81 Ash \m{4}, \textbf{75 Bh} \m{4}}
\gap$\cdot$\gap \mk{third-most-asked topic in the whole subject}
\gap$\cdot$\gap two riders recur: \emph{``clustering falls under which?''}
(\yr{\textit{74 Ma}}, \yr{\textbf{\textit{72 Ash}}}) and \emph{``learning is essential
for intelligent agents, comment''} (\yr{\textbf{\textit{74 Bh}}}, \yr{79 Jth})
\gap$\cdot$\gap \yr{78 Po} wants all \emph{three} types with examples}

\sbsr{0.60}{%
\begin{tabularx}{\linewidth}{@{}L{2.2cm}XX@{}}
\hrow \thead{} & \thead{Supervised} & \thead{Unsupervised} \\
Training data & \textbf{labelled}: every input carries its correct output &
  \textbf{unlabelled}: inputs only \\
Aim & learn the input $\rightarrow$ output mapping, then predict for new inputs &
  find hidden structure, groups and patterns in the data \\
Feedback & direct: error $=$ predicted $-$ actual, at every sample & none;
  the model is judged only by the structure it finds \\
Tasks & \textbf{classification}, \textbf{regression} &
  \textbf{clustering}, \textbf{association} \\
Algorithms & linear + logistic regression, SVM, naive Bayes, KNN, decision tree,
  back-propagation ANN & K-means, hierarchical (agglomerative), DBSCAN, Apriori, SOM \\
Human effort & high: someone has to label everything & low \\
Accuracy & higher, and measurable & lower, and hard to measure \\
Example & spam filter trained on mails already marked spam / not spam &
  segment customers by purchase behaviour, with no idea of the segments beforehand \\
\end{tabularx}}{%
\lead{Reinforcement learning}
\figT{c6_rl_loop.png}
\begin{itemize}
  \item \textbf{Feedback-based}: the agent acts, the environment returns a
        \textbf{reward} or a \textbf{penalty}, and the agent adjusts.
  \item No labelled data and no supervisor. The agent \mk{discovers}
        which actions pay by trying them.
  \item Goal $=$ maximise \textbf{cumulative} reward over the long run, not the next
        reward.
  \item \textbf{Four elements to name} (\yr{78 Po}):
  \begin{itemize}
    \item \textbf{Policy} $-$ state $\rightarrow$ action mapping.
    \item \textbf{Reward signal} $-$ one number per step, defines the goal.
    \item \textbf{Value function} $-$ total reward expected \emph{from} a state:
          good in the long run, where reward is good in the short run.
    \item \textbf{Model} of the environment (optional).
  \end{itemize}
  \item Example: a robot that gets $+10$ for reaching the goal square and $-50$ for
        falling in a pit. Also game playing, robot navigation, self-driving control.
  \item \textbf{Semi-supervised} is the fourth type worth one line: a small labelled set
        plus a large unlabelled one, used when labelling everything is too costly.
\end{itemize}}

\sbs{%
\lead{The rider: which category is clustering? \gap{\color{sub}\footnotesize(\yr{\textit{74 Ma}}, \yr{\textbf{\textit{72 Ash}}}, \m{4} of the 6)}}
\begin{itemize}
  \item \mk{Clustering is unsupervised}, because the data carries
        \mk{no class label} $-$ the algorithm both invents the groups
        and assigns points to them.
  \item \textbf{Justify by naming algorithms}, which is what the question asks for and
        where the marks are:
  \begin{itemize}
    \item \textbf{K-means} $-$ choose $k$, assign every point to the nearest centroid,
          recompute the centroids, repeat until nothing moves. Nobody ever tells it a
          right answer.
    \item \textbf{Hierarchical (agglomerative)} $-$ start with every point its own
          cluster, repeatedly merge the two nearest, stop at $k$ clusters or one.
    \item \textbf{DBSCAN} $-$ groups by density, and finds the number of clusters itself.
    \item \textbf{SOM / Kohonen} (Ch\,7) is a clustering network, and is unsupervised for
          exactly the same reason.
  \end{itemize}
  \item \mk{Contrast with classification}, which draws the same
        picture but is supervised: the classes exist \emph{before} training and every
        training point already carries one.
\end{itemize}}{%
\lead{Deciding the type of any example the examiner invents}
\begin{tabularx}{\linewidth}{@{}L{4.0cm}X@{}}
\hrow \thead{What does the data carry?} & \thead{Then it is} \\
a correct output for every input & supervised \\
nothing but inputs & unsupervised \\
a few labelled, most not & semi-supervised \\
no outputs, but a reward after acting & reinforcement \\
\end{tabularx}
\begin{itemize}
  \item Worked on the examples the papers use: spam filter $=$ supervised
        (mails already marked); customer segmentation $=$ unsupervised; game-playing
        agent $=$ reinforcement; handwriting recognition trained on labelled digits $=$
        supervised; market-basket association $=$ unsupervised.
  \item \mk{One sentence that is always worth writing}: supervised
        learning is told the answer and learns the mapping; unsupervised learning is told
        nothing and learns the structure; reinforcement learning is told only whether it
        did well, and learns the policy.
\end{itemize}}

%-----------------------------------------------------------------------
\T{6.4 Learning by analogy}

\Q{What is learning by analogy? Give an example, with a semantic network.}{\m{3}\m{2+6}\m{1+7}}
{\color{sub}\footnotesize \tP{3} \yr{\textbf{71 Bh} \m{1+7}, 76 Ba \m{2+6}, 71 Ma \m{3}}
\gap$\cdot$\gap \yr{\textbf{71 Bh}} demands the semantic network, so the answer is a
\emph{drawing plus the mapping table} \gap$\cdot$\gap \yr{71 Ma} asks it as
\emph{analogy Vs inductive learning}}

\sbsr{0.53}{%
\begin{itemize}
  \item \textbf{Learning by analogy $=$ acquiring knowledge about a new entity by
        transferring it from a known, similar entity.}
  \item It carries a solution \mk{across domains}: find the analogous
        states and operators in the two domains, then reuse the old solution.
  \item \mk{One example is enough}, which is what makes it different
        from inductive learning (which needs many).
  \item \textbf{The five steps} (write these for a 6 or 7 mark answer):
  \begin{itemize}
    \item \textbf{1. Recognition} $-$ notice the new problem resembles a past one.
    \item \textbf{2. Access and recall} $-$ the resemblance indexes memory; retrieve one
          or more candidate analogues.
    \item \textbf{3. Selection and mapping} $-$ pick the relevant parts and map them from
          the \textbf{base} (known) domain to the \textbf{target} (new) one.
    \item \textbf{4. Extension} $-$ modify and extend the mapped solution to fit the
          target.
    \item \textbf{5. Validation and generalization} $-$ test it; if it holds, form a
          general rule covering both situations.
  \end{itemize}
  \item \textbf{Analogy Vs induction} \gap{\color{sub}\footnotesize(\yr{71 Ma} \m{3})}:
        analogy maps \emph{one} case onto \emph{one} case and reasons by similarity;
        induction generalises \emph{many} cases into a rule. Analogy transfers,
        induction abstracts.
\end{itemize}}{%
\lead{First example: hydraulics and Kirchhoff}
\figT{c6_analogy.png}
\begin{itemize}
  \item Flow into a junction equals flow out; current into a node equals current out.
        Knowing $Q_a=3$, $Q_b=9$ gives $Q_c$ the same way $I_3=I_1+I_2$ does.
  \item \textbf{The mapping}: pressure drop $\leftrightarrow$ voltage drop, pipe
        $\leftrightarrow$ wire, flow rate $\leftrightarrow$ current, junction
        $\leftrightarrow$ node.
  \item So a student who knows Kirchhoff's law can solve the hydraulics problem
        \mk{without being taught hydraulics}. That sentence is the
        definition made concrete, and it is what the examiner wants.
\end{itemize}}

\sbsr{0.53}{%
\begin{itemize}
  \item \textbf{Second example: the atom and the solar system}, the one to draw as a
        \textbf{semantic net} for \yr{\textbf{71 Bh}}: two small nets, one arc list
        mapped onto the other.
  \begin{itemize}
    \item Base net: \code{Sun --greater-mass--> Earth},
          \code{Sun --attracts--> Earth}, \code{Earth --revolves-around--> Sun}.
    \item Target net: \code{Nucleus --greater-mass--> Electron},
          \code{Nucleus --attracts--> Electron}, and the arc the analogy
          \mk{supplies}: \code{Electron --revolves-around--> Nucleus}.
    \item Draw the two nets side by side and \mk{a dashed arrow
          between the matching nodes}: Sun $\leftrightarrow$ Nucleus,
          Earth $\leftrightarrow$ Electron. That dashed mapping \emph{is} the answer.
  \end{itemize}
\end{itemize}}{%
\lead{The mapping table, which is the marked half}
\begin{tabularx}{\linewidth}{@{}L{2.9cm}XX@{}}
\hrow \thead{Role} & \thead{Base: solar system} & \thead{Target: atom} \\
Heavy central body & Sun & Nucleus \\
Light orbiting body & Earth & Electron \\
Relation known & attracts, greater mass & attracts, greater mass \\
Relation \emph{transferred} & Earth revolves around Sun &
  \mk{Electron revolves around Nucleus} \\
\end{tabularx}
\begin{itemize}
  \item The last row is the whole point: it is \mk{not observed, it is
        carried across} by the analogy, and then validated.
  \item Say in one line what analogy risks: the mapping may not hold
        (electrons do not really orbit like planets), so step 5 validation is not
        optional. Naming that limitation reliably earns a mark.
\end{itemize}}

%-----------------------------------------------------------------------
\T{6.5 Inductive learning, decision trees and ID3}

\Q{Define inductive learning. How is the best attribute selected in a decision tree? Explain ID3 with an example.}{\m{8}\m{1+7}\m{2+6}\m{2+7}}
{\color{sub}\footnotesize \tF{6} \yr{\textbf{\textit{74 Bh}} \m{2+7},
\textit{74 Ma} \m{8}, \textbf{\textit{72 Ash}} \m{8}, 79 Jth \m{1+7},
\textbf{\textit{77 Ch}} \m{8}, \textbf{74 Bh} \m{2+6}} \gap$\cdot$\gap
\mk{five of the six print a data table and want the root attribute
computed}, so this section is theory only $-$ all three tables are worked in the
companion \gap$\cdot$\gap \yr{\textbf{\textit{77 Ch}}} disguises it as the
\emph{mushroom island} story}

\sbsr{0.52}{%
\lead{Inductive learning, and the tree}
\begin{itemize}
  \item \textbf{Inductive learning $=$ drawing a general hypothesis from a set of
        specific input--output examples}, chosen so that it classifies the examples well
        \emph{and} generalises to new ones.
  \item A \textbf{decision tree} is the standard hypothesis form:
  \begin{itemize}
    \item each \textbf{internal node} tests one attribute,
    \item each \textbf{branch} is one value of that attribute,
    \item each \textbf{leaf} carries a class label.
  \end{itemize}
  \item To classify, walk root to leaf and read off the label.
  \item \textbf{ID3} (Iterative Dichotomiser 3, Quinlan) builds that tree
        \mk{greedily, top down}: at every node it puts the attribute
        that separates the classes best, and never goes back.
\end{itemize}
\lead{The two formulas $-$ memorise these}
\[ H(S) \;=\; -\sum_{i} p_i \log_2 p_i \]
{\color{sub}\footnotesize entropy of set $S$: $p_i$ is the fraction of $S$ in class $i$.
$H=0$ when $S$ is pure, $H=1$ when a two-class set is split 50/50.}
\[ Gain(S,A) \;=\; H(S) \;-\; \sum_{v \in A} \frac{|S_v|}{|S|}\, H(S_v) \]
{\color{sub}\footnotesize the second term is the \emph{weighted average} entropy after
splitting on $A$. Gain is the uncertainty removed.}
\begin{itemize}
  \item \mk{Highest gain $=$ lowest weighted entropy $=$ best
        attribute.} The two rules are the same rule.
\end{itemize}}{%
\lead{Why it works, in one figure}
\figT{c6_id3_split.png}
{\color{sub}\footnotesize Splitting on \emph{Type} leaves every branch as mixed as the
parent: gain 0, useless. Splitting on \emph{Patrons} produces two pure branches and one
mixed one: high gain. That comparison \emph{is} the answer to ``how is the best attribute
selected''.}}

\sbsr{0.52}{%
\lead{The algorithm, as five steps}
\begin{itemize}
  \item \textbf{1.} Compute $H(S)$ for the current set.
  \item \textbf{2.} For each unused attribute $A$, split $S$ by its values and compute the
        weighted entropy of the split.
  \item \textbf{3.} $Gain(S,A) = H(S) -$ that weighted entropy.
  \item \textbf{4.} Make the attribute with the \textbf{largest gain} the node, and branch
        on its values.
  \item \textbf{5.} Recurse on each branch with the remaining attributes. Stop a branch
        when: it is \textbf{pure} (leaf $=$ that class); \textbf{no attributes are left}
        (leaf $=$ majority class); or it is \textbf{empty} (leaf $=$ parent's majority).
\end{itemize}
\begin{itemize}
  \item \mk{Three arithmetic checks that save marks}: $H$ of a pure
        branch is \textbf{0}, not 1; $H$ of an even two-class split is exactly
        \textbf{1}; and every gain must be $\geq 0$, so a negative one is an error in
        your weighted sum.
\end{itemize}}{%
\lead{Characteristics to list if asked}
\begin{itemize}
  \item Greedy, so \mk{no guarantee of the optimal tree} $-$ it can
        settle in a local optimum.
  \item Produces small trees, but not always the smallest.
  \item \textbf{Overfits} the training data; prefer smaller trees, and prune.
  \item Awkward on \textbf{continuous} attributes: too many possible split points.
  \item Handles noise and missing values badly. C4.5 fixes all four.
\end{itemize}}

%-----------------------------------------------------------------------
\T{6.6 Genetic algorithm}

\Q{What is a genetic algorithm? Explain all its steps with a block diagram, and its operators.}{\m{3}\m{4}\m{5}\m{7}\m{8}\m{10}\m{2+8}}
{\color{sub}\footnotesize \tS{21} \yr{69 Po \m{10}, 71 Ma \m{2+8}, 72 Ma \m{2+6+2},
75 Ba \m{8}, \textbf{77 Ch} \m{8}, \textbf{\texttt{81 Bh}} \m{2+6},
\textbf{79 Ch} \m{2+6}, \texttt{80 Ba} \m{8}, \textbf{\textit{76 Bh}} \m{5},
\texttt{81 Ba} \m{3+5}, \textit{72 Ma} \m{7}, 82 Ka \m{2+6},
\textbf{\texttt{82 Bh}} \m{2+6}, 81 Ash \m{3+5}, \textbf{81 Ch} \m{7},
\textbf{78 Ch} \m{2+6}, \textbf{\textit{74 Bh}} \m{4}, \textbf{\textit{73 Bh}} \m{3},
\textbf{70 Bh} \m{2+6}, \textbf{74 Bh} \m{4+4}, \textbf{75 Bh} \m{2+8}}
\gap$\cdot$\gap \mk{the single most-asked topic in the chapter and the
second most-asked in the subject} \gap$\cdot$\gap it has been set in every year from
2069 to 2082 \gap$\cdot$\gap the wording rotates but the answer never does: \emph{steps}
$+$ \emph{block diagram} $+$ \emph{the three operators} $+$ \emph{a worked example}}

\sbsr{0.47}{%
\lead{Definition and idea}
\begin{itemize}
  \item \textbf{A genetic algorithm is a search and optimization technique that imitates
        natural selection}: it keeps a \textbf{population} of candidate solutions and
        evolves it, generation by generation, towards better ones.
  \item Darwin in one line: \mk{select the best, discard the rest}.
  \item \textbf{Evolutionary computing} (asked by name in \yr{\textbf{79 Ch}}) is the
        family; GA is its best-known member, alongside genetic programming, evolutionary
        strategies and swarm methods.
  \item Formally a variant of \textbf{stochastic beam search}: many states kept at once,
        successors generated by combining pairs rather than mutating one.
  \item \textbf{Vocabulary} $-$ the examiner marks these:
  \begin{itemize}
    \item \textbf{Gene} $-$ one position in the string; controls one characteristic.
    \item \textbf{Chromosome} (individual) $-$ a chain of genes $=$ \textbf{one candidate
          solution}, usually a binary string.
    \item \textbf{Population} $-$ the set of chromosomes alive in one generation.
    \item \textbf{Fitness function} $-$ scores a chromosome. \mk{Problem
          dependent; it is the only thing that knows what ``good'' means.}
    \item \textbf{Generation} $-$ one full cycle of the loop.
  \end{itemize}
\end{itemize}}{%
\figT{c6_ga_flowchart.png}
{\color{sub}\footnotesize The flowchart \yr{\texttt{81 Ba}} asks for by name. Draw
\textbf{START $\rightarrow$ generate random population $\rightarrow$ evaluate fitness
$\rightarrow$ good enough? $\rightarrow$ (no) select and mate $\rightarrow$ mutate
$\rightarrow$ back to evaluate; (yes) END}.}}

\sbsr{0.55}{%
\lead{The five steps, which are also the five operators}
\begin{itemize}
  \item \textbf{1. Initialization} $-$ generate the starting population at random.
  \begin{itemize}
    \item Encode a solution as a string first (\textbf{chromosome encoding}); binary is
          usual, but integer and permutation encodings exist.
    \item Population size trades diversity against run time. Seed it around a promising
          region if you have one.
  \end{itemize}
  \item \textbf{2. Fitness evaluation} $-$ score every chromosome with the fitness
        function. Infeasible chromosomes get fitness 0.
  \item \textbf{3. Selection} $-$ choose parents, with probability rising with fitness.
  \begin{itemize}
    \item \textbf{Roulette wheel}: slice proportional to fitness.
          \textbf{Tournament}: pick $k$ at random, take the best.
          \textbf{Rank}: order them, select on rank.
          \textbf{Elitism}: copy the best one or two through untouched.
    \item \mk{Keep a few weak ones on purpose} $-$ they hold the
          diversity the next crossover needs.
  \end{itemize}
  \item \textbf{4. Crossover (recombination)} $-$ cut two parents and swap the tails, so
        each child inherits from both. Probability typically \textbf{0.6 to 0.9}.
  \item \textbf{5. Mutation} $-$ flip a random gene with a small probability, typically
        \textbf{0.001 to 0.01}.
  \item \textbf{6. Replacement, then termination} $-$ offspring replace the weakest
        members; stop on a \textbf{fixed number of generations}, a \textbf{fitness
        threshold}, or when the population has \textbf{converged} (no improvement for
        several generations).
\end{itemize}}{%
\figT{c6_ga_cycle.png}
{\color{sub}\footnotesize The same algorithm as one cycle, which is the ``block diagram''
most papers want: population $\rightarrow$ \emph{selection} $\rightarrow$ parents
$\rightarrow$ \emph{crossover, mutation} $\rightarrow$ offspring $\rightarrow$
\emph{replacement} $\rightarrow$ population.}
\lead{Crossover, drawn}
\figT{c6_ga_crossover.png}
{\color{sub}\footnotesize Single-point: one cut, tails swapped. \textbf{Two-point}: two
cuts, the middle segment swapped. \textbf{Uniform}: each gene taken from either parent by
a coin toss.}}

\sbsr{0.45}{%
\figT{c6_ga_mutation.png}
{\color{sub}\footnotesize Mutation flips one gene of one child.}
\lead{Why mutation matters \gap{\color{sub}\footnotesize(\yr{\texttt{81 Ba}} asks exactly this)}}
\begin{itemize}
  \item Crossover only \mk{recombines genes already in the
        population}. If every chromosome carries a 0 in some position, no amount
        of crossover will ever produce a 1 there.
  \item Mutation is the \mk{only source of new genetic material}. It
        keeps diversity up and stops the search converging early on a
        \textbf{local optimum}.
  \item So the rate must be \emph{small}: too high and the search degenerates into a
        random walk, losing the good solutions selection just found.
\end{itemize}}{%
\lead{Mutation Vs crossover \gap{\color{sub}\footnotesize(\yr{\textbf{\textit{73 Bh}}} \m{3})}}
\begin{tabularx}{\linewidth}{@{}L{1.6cm}XX@{}}
\hrow \thead{} & \thead{Crossover} & \thead{Mutation} \\
Parents & two & one \\
Does & exchanges segments between two chromosomes & flips a gene in one chromosome \\
Rate & high, 0.6--0.9 & low, 0.001--0.01 \\
Role & \textbf{exploitation}: combines what already works & \textbf{exploration}:
  introduces what is missing \\
If absent & population stagnates as copies of a few parents & search sticks in a local
  optimum \\
\end{tabularx}}

\sbsr{0.50}{%
\lead{When shall we use a genetic algorithm? \gap{\color{sub}\footnotesize(\yr{81 Ash}, \yr{\textbf{80 Ch}})}}
\figT{c6_ga_optima.png}
{\color{sub}\footnotesize One global and four local optima. A hill-climber stops on
whichever bump it started under; a population searches all five at once.}
\begin{itemize}
  \item The search space has \mk{many local optima} and gradient
        methods get trapped.
  \item The objective function is \textbf{not smooth}, not differentiable, or is only
        available as a black box, so no derivative method applies.
  \item The number of parameters is \textbf{very large} and exhaustive search is
        impossible.
  \item The objective is \textbf{noisy or stochastic}.
  \item A \mk{good-enough} answer is acceptable: GA gives no optimality
        guarantee.
  \item \mk{Do not use it} when an exact algorithm exists, when a single
        evaluation of the fitness function is expensive, or when the answer must be
        provably optimal.
\end{itemize}}{%
\lead{Where it is applied \gap{\color{sub}\footnotesize(\yr{72 Ma} \m{2})}}
\begin{itemize}
  \item Optimization: travelling salesman, knapsack, job-shop scheduling, timetabling.
  \item Engineering design, antenna and circuit layout, PID controller tuning.
  \item Neural network weight and architecture search; feature selection in ML.
  \item Robotics path planning, vehicle routing, bioinformatics (gene sequencing).
  \item \textbf{Why the study of genes belongs in AI} (\yr{\textbf{74 Bh}} \m{4}):
        evolution is a \mk{proven search procedure} that solved design
        problems no designer could; GA borrows the mechanism, not the biology.
\end{itemize}
\lead{Advantages and limitations}
\begin{itemize}
  \item \checkmark Simple concept; works on the \textbf{encoded} parameters, so it needs
        no derivative and no auxiliary information, only a fitness score.
  \item \checkmark Searches many points at once, so it escapes local optima; inherently
        \textbf{parallel}; the answer improves with run time and can be stopped anywhere.
  \item \xmark No guarantee of the global optimum; \xmark many fitness
        evaluations, so slow when each is costly; \xmark the encoding, population size
        and rates all need tuning; \xmark can converge prematurely if diversity is lost.
\end{itemize}}

%-----------------------------------------------------------------------
\T{6.7 Fuzzy logic and the fuzzy inference system}

\Q{What is fuzzy logic? When do we use it? Explain fuzzy inference / the steps, with an example.}{\m{2}\m{3}\m{4}\m{8}\m{2+6}\m{3+7}\m{1+2+5}}
{\color{sub}\footnotesize \tS{16} \yr{\textbf{69 Bh} \m{4}, 70 Ma \m{2+6},
80 Ash \m{1+2+5}, 73 Ma \m{3+7}, 79 Ash \m{1+2+5}, \textbf{80 Ch} \m{2+5},
\texttt{82 Ba} \m{2+6}, \textbf{\texttt{80 Bh}} \m{8}, \textbf{\texttt{79 Bh}} \m{2+6},
\texttt{81 Ba} \m{4}, 78 Po \m{4}, \textbf{77 Ch} \m{4}, 81 Ash \m{4},
\texttt{80 Ba} \m{3}, 82 Ka \m{2}, \textbf{\texttt{82 Bh}} \m{2}}
\gap$\cdot$\gap \mk{the syllabus calls this ``fuzzy learning'' and no
paper ever has}: they ask fuzzy \emph{logic} and fuzzy \emph{inference}, which is
reasoning, not learning \gap$\cdot$\gap \yr{73 Ma} is the only one to name
\textbf{Mamdani} \gap$\cdot$\gap \yr{82 Ka} and \yr{\textbf{\texttt{82 Bh}}} bolt it onto
a GA question}

\sbsr{0.50}{%
\lead{What it is}
\begin{itemize}
  \item \textbf{Fuzzy logic is a many-valued logic in which a proposition has a degree of
        truth anywhere in $[0,1]$}, instead of only 0 or 1.
  \item It exists because human concepts have \mk{no sharp boundary}:
        \emph{tall}, \emph{hot}, \emph{slow}, \emph{near}, \emph{honest}. A crisp rule
        must answer \emph{is 179\,cm tall?} yes or no; a fuzzy one answers 0.78.
  \item Introduced by \textbf{Lotfi A. Zadeh}, UC Berkeley, \textbf{1965}.
  \item \textbf{Fuzzy set}: for a universe $U$, a fuzzy set $A$ is the set of pairs
        $\{(x, \mu_A(x)) : x \in U\}$ where the \textbf{membership function}
  \[ \mu_A(x): U \rightarrow [0,1] \]
  \begin{itemize}
    \item $\mu_A(x)=1$ $\rightarrow$ $x$ is fully in $A$;
          $\mu_A(x)=0$ $\rightarrow$ fully out;
          $0<\mu_A(x)<1$ $\rightarrow$ partly in.
  \end{itemize}
  \item \textbf{Linguistic variable} $-$ a variable whose values are words:
        Temperature $=$ \{very cold, cold, warm, hot, very hot\}. Each word is a fuzzy
        set with its own membership curve, and the curves \textbf{overlap} on purpose.
\end{itemize}}{%
\figT{c6_crisp_fuzzy_tbl.png}
{\color{sub}\footnotesize \textbf{Crisp Vs fuzzy on the same data.} Membership of
\emph{tall}: the crisp column steps from 0 to 1 at 180\,cm, so Tom at 181 is tall and
David at 179 is not. The fuzzy column grades them 0.82 and 0.78 $-$ which is what a person
would say. This one table answers \emph{``how does fuzzy logic differ from classical
binary logic''} (\yr{\textbf{\texttt{82 Bh}}}) outright.}
\figT{c6_crisp_fuzzy_set.png}
{\color{sub}\footnotesize The same point as sets: $b$ is simply outside crisp $A$; in
fuzzy $A$ the boundary is a band, so $b$ and $c$ have partial membership.}}

\sbsr{0.48}{%
\lead{Fuzzy set operations}
{\color{sub}\footnotesize all defined pointwise on the membership values; worked on real
numbers in the companion.}
\begin{tabularx}{\linewidth}{@{}L{3.1cm}X@{}}
\hrow \thead{Operation} & \thead{Membership} \\
Union $A \cup B$ & $\max(\mu_A, \mu_B)$ \\
Intersection $A \cap B$ & $\min(\mu_A, \mu_B)$ \\
Complement $A'$ & $1 - \mu_A$ \\
Bold union (X-OR) & $\min(1,\; \mu_A + \mu_B)$ \\
Bold intersection & $\max(0,\; \mu_A + \mu_B - 1)$ \\
Equality & $\mu_A(x) = \mu_B(x)$ for every $x$ \\
\end{tabularx}
\lead{Fuzzy logic Vs probability}
\begin{itemize}
  \item \textbf{Not the same thing, and papers reward saying so.}
  \item \emph{Ram's membership of ``old'' is 0.9} $-$ Ram \mk{is} 90\,\%
        old. A statement about \textbf{vagueness}, and it stays 0.9 after you meet him.
  \item \emph{P(Ram is old) $=$ 0.9} $-$ Ram either is or is not old and we do not know
        which. A statement about \textbf{ignorance}, and it collapses to 0 or 1 once you
        meet him.
\end{itemize}}{%
\lead{Architecture of a fuzzy inference system}
\figT{c6_fuzzy_arch.png}
{\color{sub}\footnotesize The block diagram \yr{\texttt{82 Ba}} and five other papers ask
for: crisp input $\rightarrow$ \textbf{fuzzifier} $\rightarrow$ fuzzy input set
$\rightarrow$ \textbf{inference engine} (driven by the \textbf{rule base} above it)
$\rightarrow$ fuzzy output set $\rightarrow$ \textbf{defuzzifier} $\rightarrow$ crisp
output.}
\begin{itemize}
  \item \textbf{Rule base} $-$ the IF--THEN rules supplied by domain experts.
        \emph{IF temperature is cold AND target is warm THEN heater is on.}
  \item \textbf{Fuzzification} $-$ read the crisp sensor value off each membership curve
        to get a degree for every linguistic term.
  \item \textbf{Inference engine} $-$ match the fuzzified inputs against the rules,
        compute each rule's \textbf{firing strength}, and combine the consequents.
  \item \textbf{Defuzzification} $-$ collapse the aggregated fuzzy output back to one
        crisp number. \textbf{Centroid} (centre of gravity) is the usual method; also
        mean of maxima, bisector, weighted average.
\end{itemize}}

\sbsr{0.52}{%
\lead{The steps of ``fuzzy learning'' \gap{\color{sub}\footnotesize(\yr{\textbf{\texttt{80 Bh}}} \m{8}, \yr{\textbf{\texttt{79 Bh}}} \m{2+6}, \yr{\textbf{69 Bh}} \m{4})}}
\begin{itemize}
  \item \textbf{1. Define the linguistic variables and terms.}
        Temperature $=$ \{very cold, cold, warm, very warm, hot\}.
  \item \textbf{2. Construct a membership function for each term.} Triangular and
        trapezoidal are normal; they must overlap.
  \item \textbf{3. Build the rule base} of IF--THEN rules, usually as a matrix of every
        input combination against the action to take.
  \item \textbf{4. Fuzzify}: convert the crisp readings into membership degrees.
  \item \textbf{5. Evaluate the rules} in the inference engine, and \textbf{aggregate}
        the results into one fuzzy output.
  \item \textbf{6. Defuzzify} into a crisp control value.
  \item \mk{Worked end to end on an air-conditioner in the companion.}
\end{itemize}}{%
\figT{c6_memb_fns.png}
{\color{sub}\footnotesize Step 2 for three linguistic variables of a soil-moisture
controller. Note the overlaps: at 45 the humidity is \emph{both} medium and wet, to some
degree.}}

\sbsr{0.52}{%
\lead{Mamdani fuzzy inference \gap{\color{sub}\footnotesize(\yr{73 Ma} \m{3+7})}}
\begin{itemize}
  \item Built by \textbf{Ebrahim Mamdani}, London, \textbf{1975}, to control a steam
        engine and boiler from rules given by human operators.
  \item Its six operations, in order: \textbf{determine the rules}, \textbf{fuzzify} the
        inputs, combine them to get each rule's \textbf{strength} (fuzzy AND $=$ min),
        apply \textbf{implication} to clip the output set, \textbf{aggregate} the clipped
        sets, then \textbf{defuzzify}.
  \item \mk{Mamdani Vs Sugeno} in one line: Mamdani's rule consequent is
        a \emph{fuzzy set} and needs defuzzifying; Sugeno's is a \emph{function of the
        inputs}, so the output is a weighted average and no defuzzification is needed.
        Mamdani is more intuitive, Sugeno cheaper.
\end{itemize}}{%
\figT{c6_fuzzy_infer.png}
{\color{sub}\footnotesize The same architecture with the membership curves drawn on it, so
``fuzzy input'' and ``fuzzy output'' are visible as shapes rather than words.}}

\sbsr{0.50}{%
\lead{Advantages}
\begin{itemize}
  \item Mimics human reasoning, so the rules are \mk{readable and can be
        written by a domain expert} who knows no mathematics.
  \item Needs no exact mathematical model of the plant $-$ decisive where one does not
        exist.
  \item \textbf{Robust}: imprecise, noisy or partly missing inputs still give a sensible
        output, so it degrades gracefully when a sensor fails.
  \item Simple, cheap to implement on small hardware; handles nonlinearity of arbitrary
        complexity.
  \item Proven in real products: washing machines, air conditioners, camera autofocus,
        anti-lock brakes, subway train braking.
\end{itemize}}{%
\lead{Disadvantages, and when not to use it}
\begin{itemize}
  \item \mk{No systematic design method.} Rules and membership
        functions come from expert opinion, so two designers get two systems.
  \item Accuracy depends entirely on the quality of that expert knowledge.
  \item Validation and verification need extensive hardware testing; there is no proof of
        stability the way there is for a classical controller.
  \item Not suited to problems needing high precision, nor to pattern recognition where a
        neural network learns the mapping from data.
  \item \textbf{Do not use it} when the input space does not map naturally to an output
        space, when a crisp model already works, or where common sense or a simple
        threshold would do.
\end{itemize}}

%-----------------------------------------------------------------------
\T{6.8 Boltzmann machine}

\Q{Define the Boltzmann machine. Explain it with a suitable algorithm.}{\m{4}\m{1+7}\m{2+8}}
{\color{sub}\footnotesize \tP{3} \yr{70 Ma \m{1+7}, \textbf{72 Ash} \m{2+8}, 76 Ba \m{4}}
\gap$\cdot$\gap only three papers, but two of them put 7 or 8 marks on it
\gap$\cdot$\gap \yr{70 Ma} pairs it with the semantic network (Ch\,5), \yr{\textbf{72 Ash}}
with ``what is machine learning''}

\sbsr{0.52}{%
\begin{itemize}
  \item \textbf{A Boltzmann machine is a stochastic, energy-based, recurrent neural
        network of symmetrically connected binary units that switch on and off at
        random, with a probability set by the network's energy.}
  \item Invented by \textbf{Hinton and Sejnowski, 1985}. Named for the
        \textbf{Boltzmann distribution} it samples from.
  \item \textbf{Structure}:
  \begin{itemize}
    \item \textbf{Visible units} $V$ carry the data in \emph{and} the reconstruction out.
    \item \textbf{Hidden units} $H$ capture features and dependencies.
    \item Every pair is \textbf{symmetrically} connected, $w_{ij}=w_{ji}$, and there are
          \mk{no self connections}.
    \item \mk{There is no output layer} $-$ the same visible nodes
          that take the input return the reconstruction.
  \end{itemize}
  \item \textbf{State and energy}. Each unit is 0 or 1; the network's energy is
  \[ E \;=\; -\sum_{i<j} w_{ij}\, s_i\, s_j \;-\; \sum_i \theta_i s_i \]
  and the probability of a state falls exponentially with its energy:
  \[ p_i \;=\; \frac{e^{-\varepsilon_i / kT}}{\sum_{j} e^{-\varepsilon_j / kT}} \]
  \item So the machine \mk{prefers low-energy states but can still
        climb out of one}, which is how it escapes local minima $-$ the same
        idea as simulated annealing (Ch\,3), with $T$ lowered as it runs.
\end{itemize}}{%
\figT{c6_boltzmann.png}
{\color{sub}\footnotesize Visible nodes below, hidden above, every pair connected both
ways. Compare with a feed-forward net: no layers, no direction, no output side.}}

\sbsr{0.52}{%
\lead{The learning algorithm, in six steps}
\begin{itemize}
  \item \textbf{1.} Initialise the weights randomly and set a high temperature $T$.
  \item \textbf{2. Positive phase (clamped)}: fix the visible units to a training vector,
        let the hidden units settle, record the correlation
        $\langle s_i s_j \rangle_{data}$.
  \item \textbf{3. Negative phase (free)}: release every unit, let the network run to
        equilibrium, record $\langle s_i s_j \rangle_{model}$.
  \item \textbf{4. Update}
        $\Delta w_{ij} = \eta \left( \langle s_i s_j \rangle_{data} -
        \langle s_i s_j \rangle_{model} \right)$.
  \item \textbf{5.} Lower $T$ a little (annealing) and repeat over the training set.
  \item \textbf{6.} Stop when the weights stop changing.
\end{itemize}
\begin{itemize}
  \item \mk{What the two phases mean}: the clamped phase measures how
        units co-fire when the world is showing them the truth, the free phase how they
        co-fire when the network is dreaming on its own. Learning pushes the second
        towards the first.
\end{itemize}}{%
\lead{Restricted Boltzmann machine (RBM)}
\begin{itemize}
  \item A full Boltzmann machine is \mk{far too slow}: every unit
        depends on every other, so the weights must be updated one at a time.
  \item RBM \textbf{forbids connections inside a layer}. Visible-to-hidden only.
  \item Units within a layer are then independent given the other layer, so a whole layer
        updates \textbf{in parallel}. Trained by \textbf{contrastive divergence}.
  \item Used for dimensionality reduction, classification, regression, collaborative
        filtering (recommender systems), and as the building block of deep belief
        networks.
\end{itemize}
\figT{c6_rbm.png}}

\hr

%-----------------------------------------------------------------------
\T{6.9 Every Chapter 6 question, grouped by the answer it wants}

{\color{sub}\footnotesize Every Chapter 6 question in the 41 papers, grouped by the answer
it wants. \textbf{Bold} $=$ regular exam, plain $=$ back. \texttt{Monospace} $=$ CT 710
(BEI), \textit{italic} $=$ CT 78506 (BEX), upright $=$ CT 653 (BCT).}

\begin{itemize}
  \item \textbf{Genetic algorithm} \textbf{21 papers} $\rightarrow$ \S6.6, worked in the companion.
        The wordings, all one answer: \emph{explain GA and its operators}
        (\yr{69 Po}, \yr{71 Ma}, \yr{72 Ma}); \emph{all steps with block diagram}
        (\yr{75 Ba}, \yr{\textbf{77 Ch}}, \yr{\textbf{\texttt{81 Bh}}},
        \yr{\texttt{80 Ba}}, \yr{\textbf{\textit{76 Bh}}}, \yr{\textbf{78 Ch}});
        \emph{steps with a flowchart, and why mutation matters} (\yr{\texttt{81 Ba}});
        \emph{working principle with examples} (\yr{\textit{72 Ma}},
        \yr{\textbf{81 Ch}}); \emph{when shall we use it} (\yr{81 Ash});
        \emph{what is evolutionary computing} (\yr{\textbf{79 Ch}});
        \emph{the algorithm of GA} (\yr{\textbf{70 Bh}}); \emph{justify that the study of
        genes is part of AI} (\yr{\textbf{74 Bh}}); short notes
        (\yr{\textbf{\textit{74 Bh}}}, \yr{\textbf{\textit{73 Bh}}} as
        \emph{mutation Vs crossover}); and bolted to a fuzzy question in
        \yr{82 Ka} and \yr{\textbf{\texttt{82 Bh}}}.
  \item \textbf{Fuzzy logic and fuzzy inference} \textbf{16 papers} $\rightarrow$ \S6.7.
        \emph{Fuzzy inference with an example} (\yr{70 Ma}, \yr{80 Ash});
        \emph{Mamdani method} (\yr{73 Ma}); \emph{when do we use it, and the steps}
        (\yr{79 Ash}, \yr{\textbf{80 Ch}}); \emph{architecture and working mechanism}
        (\yr{\texttt{82 Ba}}); \emph{all the steps in fuzzy learning with a block diagram}
        (\yr{\textbf{\texttt{80 Bh}}}, \yr{\textbf{\texttt{79 Bh}}}, which also asks
        \emph{what is membership of an element}); \emph{what is fuzzy learning, with a
        practical example} (\yr{\textbf{69 Bh}}); and short notes at
        \yr{\texttt{81 Ba}}, \yr{78 Po}, \yr{\textbf{77 Ch}}, \yr{81 Ash},
        \yr{\texttt{80 Ba}}.
  \item \textbf{Supervised Vs unsupervised, and the three types} \textbf{12 papers}
        $\rightarrow$ \S6.3. The \emph{clustering} rider is in
        \yr{\textit{74 Ma}} and \yr{\textbf{\textit{72 Ash}}}; the
        \emph{``learning is essential for intelligent agents''} rider in
        \yr{\textbf{\textit{74 Bh}}} and \yr{79 Jth} (answered in Ch\,1); all three types
        with examples in \yr{78 Po}.
  \item \textbf{ID3 and the best attribute} \textbf{6 papers} $\rightarrow$ \S6.5 for the method,
        companion for all three data tables: \emph{play golf}
        (\yr{\textbf{\textit{74 Bh}}}, \yr{\textit{74 Ma}}, \yr{\textbf{\textit{72 Ash}}}),
        \emph{play cricket} (\yr{79 Jth}), \emph{poisonous mushrooms}
        (\yr{\textbf{\textit{77 Ch}}}). \yr{\textbf{74 Bh}} asks the ID3 \emph{process}
        without a table.
  \item \textbf{Ways of learning} \textbf{2 papers} $\rightarrow$ \S6.2
        (\yr{\textbf{\textit{71 Bh}}} \m{8}, \yr{\texttt{81 Ba}} \m{3}). \emph{Induction
        Vs deduction} as a short note at \yr{\textbf{\textit{73 Bh}}} \m{3}.
        \emph{Algorithms inspired by biology and social behaviour} at
        \yr{\textbf{78 Ch}} \m{2}.
  \item \textbf{Learning framework} \textbf{2 papers} $\rightarrow$ \S6.1
        (\yr{78 Ba}, \yr{\textbf{69 Bh}}). \emph{Define machine learning} alone at
        \yr{80 Ash} \m{1}, and as the lead-in to eleven other questions.
  \item \textbf{Learning by analogy} \textbf{3 papers} $\rightarrow$ \S6.4.
        \yr{\textbf{71 Bh}} wants it \emph{with a semantic network};
        \yr{71 Ma} wants \emph{analogy Vs inductive learning} \m{3}.
  \item \textbf{Boltzmann machine} \textbf{3 papers} $\rightarrow$ \S6.8
        (\yr{70 Ma}, \yr{\textbf{72 Ash}}, \yr{76 Ba}).
  \item \textbf{Swarm intelligence} \m{2} \added{} $-$ asked once,
        \yr{\textbf{75 Bh}}, and \mk{no lecture deck in this subject
        mentions it}. Answer: a family of optimization methods in which many
        simple agents following local rules produce intelligent group behaviour with no
        central control $-$ ant colony optimization (pheromone trails, shortest path),
        particle swarm optimization (each particle steers towards its own best and the
        swarm's best), bee colony, flocking. Like GA it is population based and
        nature inspired; unlike GA the individuals \mk{are not replaced,
        they move}. The other half of that question, the GA fitness numerical, is
        worked in the companion.
  \item \textbf{Reinforcement learning in detail} $-$ named by \yr{78 Po} only, and the
        one lecture deck that covers it is PS Sir's. \S6.3 carries the four elements
        (policy, reward signal, value function, model).
  \item \textbf{Not in this chapter}: \emph{neural networks} appear on the syllabus line
        for Chapter 6 but every paper sets them under Applications of AI $-$ perceptron,
        back-propagation, Hopfield, Kohonen and Hebbian learning are all in
        \textbf{Ch\,7}. \emph{Learning in a neural network} short notes (\yr{69 Po},
        \yr{\textbf{74 Bh}}, \yr{78 Ba}, \yr{\textbf{\textit{77 Ch}}}, \yr{81 Ash}) are
        answered there.
\end{itemize}

\lead{Source notes}
\begin{itemize}
  \item Spine: BA Sir \emph{CH-06 Machine Learning}, PS Sir \emph{6.\,Machine Learning},
        SG Sir \emph{AIChapter\_6}, and \emph{Insights on Artificial Intelligence} ch\,6
        for the worked method.
  \item \mk{Two corrections carried into the companion.} \emph{Insights}
        p167 prints $\mu_{A'}(30)=0.5$ for a complement whose correct value is
        \textbf{0.6}, and p166 labels a fuzzified pair \emph{(10, 0.3)} when the element
        is 9, not 10. Both are shown correctly there, with a note.
  \item \mk{Reinforcement learning and Mamdani inference are in PS Sir's
        deck only}, which is labelled \emph{old syllabus}. Both are on the current
        CT 710 syllabus by way of the papers, so both are kept, checked against the
        standard sources.
  \item Swarm intelligence is \added{} $-$ researched and written for
        \yr{\textbf{75 Bh}}, because no deck in \emph{Notes\textbackslash} contains the
        word.
\end{itemize}
